\section{Estimation du facteur spread}
\label{sec:spread}

Il n'existe pas de marché liquide d'options sur OAT d'où tirer la volatilité du spread. Il faut une autre
source d'information, et vérifier qu'elle identifie les paramètres $(a_y, \sigma_y, \rho)$. Le spread OAT–€STR
mêle crédit, liquidité et effets d'offre de titres et de collatéral (rareté des titres d'État remis en garantie,
programme d'émission du Trésor). \citet{MonfortRenne2014} identifient la liquidité par l'écart entre les titres
de la KfW, garantis par l'État fédéral allemand, et le Bund ; pour pricer une option sur OAT, seul le spread
total compte, et on ne le décompose pas.

\subsection{La courbe OAT n'identifie pas le facteur spread}
\label{sec:h1}

\citet[éq.~56-57]{Russo} calibrent le facteur spread sur la courbe des zéro-coupons risqués. Pour chaque
$\rho$ d'une grille sur $[-1,1]$, on minimise en $(a_y,\sigma_y)$ l'erreur relative entre les prix de marché
$\bar P^M(0,T_i)$ et les prix $\bar P^\ast(0,T_i) = e^{-\int_0^{T_i}\bar\alpha}$, puis on retient le $\rho$ de
plus faible erreur. Or le recalage~\eqref{eq:recalage} impose
\begin{equation}
\int_0^T\bar\alpha = -\log\bar P^M(0,T) + \tfrac12\bar V(0,T)\quad\Longrightarrow\quad
\bar P^\ast(0,T) = \bar P^M(0,T)\,e^{-\frac12\bar V(0,T)}.
\end{equation}
L'erreur à minimiser est donc $\big|e^{-\frac12\bar V(0,T)} - 1\big|$ : la procédure minimise la convexité
totale $\bar V = V_x + V_y + V_{xy}$ de la courbe OAT, et non un écart au marché. $V_x$ étant fixé par le facteur
taux, l'optimum pousse le terme croisé à compenser $V_x$, vers $\rho < 0$ et vers des $(a_y,\sigma_y)$ qui
rendent $V_y + V_{xy}$ aussi négatif que possible. La solution exacte est $\rho = -1$, $a_y = a_x$,
$\sigma_y = \sigma_x$ : le facteur spread annule alors le facteur taux ($x + y \equiv 0$), le taux OAT devient
déterministe et $\bar V \equiv 0$. Le tableau~\ref{tab:h1} le vérifie numériquement.

\begin{table}[htbp]
\centering\small
\caption{Procédure de calibration de Russo et al. appliquée à la courbe OAT ($a_x$ et $\sigma_x$ fixés)}
\label{tab:h1}
\setlength{\tabcolsep}{4pt}
\begin{tabular}{@{}lccccccccccc@{}}
\toprule
$\rho$ & $-1$ & $-0{,}8$ & $-0{,}6$ & $-0{,}4$ & $-0{,}2$ & $0$ & $0{,}2$ & $0{,}4$ & $0{,}6$ & $0{,}8$ & $1$ \\
\midrule
$a_y$ & 0,011 & 0,011 & 0,011 & 0,011 & 0,070 & 1,373 & 2,000 & 2,000 & 2,000 & 2,000 & 2,000 \\
$\sigma_y$ (\%) & 0,74 & 0,59 & 0,44 & 0,29 & 0,23 & 0,01 & 0,01 & 0,01 & 0,01 & 0,01 & 0,01 \\
Erreur (\%) & 0,000 & 2,788 & 4,860 & 6,291 & 7,136 & 7,408 & 7,409 & 7,410 & 7,411 & 7,412 & 7,414 \\
\bottomrule
\end{tabular}
\end{table}

L'erreur s'annule en $\rho = -1$ avec $(a_y, \sigma_y) = (a_x, \sigma_x)$. Pour $\rho \ge 0$, le facteur spread ne
peut qu'ajouter de la convexité : l'optimiseur ramène $\sigma_y$ sur sa borne inférieure, et l'erreur
résiduelle, 7,41~\%, est celle du seul terme de convexité du facteur taux, $|e^{-V_x/2} - 1|$. Appliquée à une courbe recalée par construction, la
procédure ne peut donc identifier aucun des trois paramètres. C'est le principe des modèles à shift
déterministe : la courbe initiale fixe le shift, et les paramètres de diffusion se calibrent sur des options
\citep[CM2, p.~30-31]{Hillairet2026}. Dans un modèle affine gaussien, la volatilité n'entre dans la coupe des
rendements que par le terme de convexité de la constante \citep[cours~3, éq.~II.8]{MPR2025}, et le recalage
absorbe ce terme. \citet[section~22.7.2]{BrigoMercurio2006} font le même constat pour l'intensité de défaut,
dont ils fixent la corrélation par estimation historique ; ils notent aussi qu'un G2++ calibré avec $\rho$
proche de $-1$ dégénère en modèle à un facteur (section~4.2.7), ce qui est l'optimum exact de la procédure.
\emph{L'hypothèse~H1 est rejetée.}

\subsection{L'information des séries historiques}

Dans le modèle, le spread de taux zéro de tenor $\tau$ vaut $c_\tau(t) + y(t)\,H_{a_y}(\tau)/\tau$, où $c_\tau$ est
déterministe. Sa volatilité instantanée est $\sigma_y H_{a_y}(\tau)/\tau$, et sa corrélation avec le taux €STR de
même tenor est $\rho$. Ce sont des propriétés des trajectoires, qui relèvent de la variation quadratique : elles
sont les mêmes sous la mesure historique et sous la mesure de pricing, car le changement de mesure de Girsanov ne
modifie que les dérives \citep[CM3, p.~5-6]{Hillairet2026}. En temps discret, un facteur d'actualisation
exponentiel-affine change de même la constante et la matrice autorégressive d'un VAR gaussien, jamais sa variance
\citep[cours~3, prop.~I.3, slides~24-25]{MPR2025}. Dans le modèle, l'égalité des volatilités sous les deux mesures
est donc un théorème ; appliquée aux données, c'est une hypothèse, celle d'une prime de risque de variance nulle,
qu'un facteur d'actualisation exponentiel-quadratique lèverait (cours~6, slides~116-117) et que l'absence d'options
sur OAT ne permet pas de tester. Le cadre de crédit va dans le même sens : quand le défaut n'est pas pricé,
l'intensité est la même fonction des facteurs sous $\P$ et sous $\Q$, et seule la dynamique des facteurs change
(cours~5, prop.~II.3).

Le coefficient $H_{a_y}(\tau)/\tau$ vient du prix d'un zéro-coupon : la vitesse de retour qui y figure est la
vitesse risque-neutre (cours~3, prop.~I.1, slide~17). La structure par terme des volatilités du spread identifie
donc $(a_y, \sigma_y)$ sous la mesure de pricing. Une volatilité décroissante avec le tenor signalerait un retour à
la moyenne ; une volatilité plate impose $a_y \approx 0$. Une régression sur le niveau du spread ne donnerait, elle,
qu'une vitesse de retour historique.

On utilise les séries quotidiennes du spread entre le rendement OAT générique et le taux de swap €STR de même
tenor, les rendements de pair servant d'approximation des taux zéro. Le tenor de deux ans est exclu : la série
générique OAT à deux ans saute de $-65$~pb le 22 janvier 2024 sans revenir, au changement du titre de référence, et
le court terme est dominé par les BTF et la politique monétaire. Le tableau~\ref{tab:vol_spread} donne les moments
par tenor sur trois fenêtres.

\begin{table}[htbp]
\centering\small
\caption{Volatilité annualisée du spread (pb par an) et corrélation de ses variations quotidiennes avec celles
du swap €STR de même tenor}
\label{tab:vol_spread}
\begin{tabular}{@{}lcccccc@{}}
\toprule
 & \multicolumn{2}{c}{2021-2026} & \multicolumn{2}{c}{Trois ans} & \multicolumn{2}{c}{Un an} \\
\cmidrule(lr){2-3}\cmidrule(lr){4-5}\cmidrule(lr){6-7}
Tenor & Volatilité & $\rho$ & Volatilité & $\rho$ & Volatilité & $\rho$ \\
\midrule
5 ans  & 33,9 & 0,064 & 28,8 & 0,067 & 26,5 & 0,272 \\
10 ans & 32,7 & 0,162 & 30,5 & 0,182 & 31,9 & 0,434 \\
15 ans & 33,2 & 0,118 & 29,8 & 0,147 & 28,2 & 0,387 \\
20 ans & 33,1 & 0,152 & 30,5 & 0,181 & 30,9 & 0,373 \\
25 ans & 33,1 & 0,148 & 30,9 & 0,169 & 30,1 & 0,360 \\
30 ans & 34,2 & 0,124 & 32,2 & 0,146 & 31,9 & 0,275 \\
\midrule
Ajustement & \multicolumn{2}{c}{$\sigma_y = 0{,}337\,\%$} & \multicolumn{2}{c}{$\sigma_y = 0{,}307\,\%$} & \multicolumn{2}{c}{$\sigma_y = 0{,}302\,\%$} \\
\bottomrule
\end{tabular}
\end{table}

La volatilité du spread est presque plate en tenor. Ajustée fenêtre par fenêtre, la structure par terme
$\sigma_y H_{a_y}(\tau)/\tau$ place $a_y$ sur sa borne inférieure de 0,001, soit une demi-vie de près de sept
siècles, dans les trois cas, avec un écart maximal de 1,0 à 3,6~pb (figure~\ref{fig:spread_historique}). La
corrélation est positive : le spread s'élargit quand les taux montent. Elle varie toutefois de 0,064 à 0,434
selon le tenor et la fenêtre, et cette fourchette servira de sensibilité.

Le modèle ne comporte qu'un facteur de spread. L'analyse en composantes principales des variations quotidiennes
du spread aux tenors 5 à 30 ans attribue 85,8~\% de leur variance à la première composante, dont les chargements
sont presque égaux d'un tenor à l'autre (de 0,36 à 0,43). Le spread se déplace surtout en bloc, ce que décrit
un facteur unique à vitesse de retour faible. La deuxième composante, une pente qui explique 8,4~\% de la variance,
reste hors du modèle.

\begin{figure}[htbp]
\centering
\includegraphics[width=\textwidth]{spread_historique}
\caption{À gauche, volatilité du spread OAT–€STR par tenor sur trois fenêtres (points) et structure par terme
ajustée (traits), comparée au profil qu'imposerait $a_y = 0{,}2$ ; à droite, variations quotidiennes du spread et
du swap €STR à vingt ans.}
\label{fig:spread_historique}
\end{figure}

\subsection{Estimation par la méthode des moments généralisée}

Le pricing n'a besoin que de $(a_y,\sigma_y,\rho)$, pas de toute la dynamique historique du spread : c'est le
cadre de la méthode des moments généralisée \citep[cours~5, slide~8]{Hansen1982,MPR2025}. On retient
douze moments, calculés sur 1\,451 variations quotidiennes aux dates communes des tenors 5, 10, 15, 20, 25 et
30 ans :
\begin{equation}
m^{vol}_\tau(\theta) = \sigma_y\,\frac{H_{a_y}(\tau)}{\tau},\qquad m^{corr}_\tau(\theta) = \rho,
\end{equation}
le premier pour la volatilité annualisée du spread, le second pour sa corrélation avec le swap €STR de même tenor.
Les corrélations entre spreads de tenors différents sont écartées : égales à 1 dans un modèle à un facteur,
elles le rejetteraient par construction.

Avec $\hat m$ les moments observés et $g(\theta) = m(\theta) - \hat m$, l'estimateur minimise
$g(\theta)'\,W\,g(\theta)$. La covariance $S$ des moments estimés doit tenir compte de deux propriétés des
données. Les variations quotidiennes du spread sont hétéroscédastiques : le test ARCH-LM d'\citet{Engle1982}, à
dix retards, rejette l'homoscédasticité à chaque tenor (statistique de 101 à 161, p-valeur au plus
$3{,}2\times10^{-17}$). Elles sont aussi légèrement autocorrélées. On estime donc $S$ par bootstrap par blocs
mobiles \citep{Kunsch1989}, avec des blocs de vingt jours consécutifs et 500 tirages, par le même code pour
tous les moments.

La pondération asymptotiquement efficace est $W = S^{-1}$. Mais les erreurs d'estimation des moments sont ici
très corrélées d'un tenor à l'autre, de 0,90 à 0,98 pour les volatilités et de 0,82 à 0,98 pour les corrélations
entre tenors voisins, et $S^{-1}$ combine alors les moments avec des poids implicites négatifs. L'estimation
efficace de $\rho$, 0,028, sort ainsi de l'intervalle de toutes les corrélations observées sur la période
(de 0,064 à 0,162). C'est le biais de petit échantillon de la GMM efficace appliquée aux structures de covariance,
documenté par \citet{AltonjiSegal1996}, qui recommandent une pondération diagonale. On retient donc
$W_1 = \mathrm{diag}(S)^{-1}$ : chaque moment pèse en proportion inverse de sa variance. Les écarts-types sont ceux
de la formule sandwich $(G'W_1G)^{-1}G'W_1SW_1G(G'W_1G)^{-1}$, où $G$ est la jacobienne de $g$, valide quelle que
soit la pondération. La GMM efficace sert au test de sur-identification de \citet{Hansen1982}. Comme $a_y$ bute sur
sa borne, l'inférence classique ne s'applique pas à ce paramètre : écarts-types et statistique $J$ sont calculés à
$a_y$ fixé.

\begin{table}[htbp]
\centering\small
\caption{Estimation GMM des paramètres du spread (écarts-types entre parenthèses)}
\label{tab:gmm}
\begin{tabular}{@{}lccc@{}}
\toprule
Pondération & $a_y$ & $\sigma_y$ (\%) & $\rho$ \\
\midrule
Diagonale, $W_1 = \mathrm{diag}(S)^{-1}$ (retenue) & 0,001 (borne) & 0,3374 (0,0182) & 0,130 (0,040) \\
Efficace, $W_2 = S^{-1}$ & 0,001 (borne) & 0,3257 (0,0141) & 0,028 (0,034) \\
\midrule
\multicolumn{4}{@{}l}{Test de Hansen : $J = 73{,}7$ pour 10 degrés de liberté, p-valeur $8{,}7\times10^{-12}$} \\
\bottomrule
\end{tabular}
\end{table}

L'estimation retenue (tableau~\ref{tab:gmm}) donne une volatilité du spread de 33,7~pb par an et une corrélation de
0,13, la vitesse de retour restant sur sa borne. Les volatilités ajustées vont de 33,2 à 33,7~pb selon le tenor,
pour des volatilités observées de 32,7 à 34,2~pb. Le test $J$ rejette le modèle. Avec près de 1\,450 observations,
le test est très puissant : il rejette un spread à un facteur qui ne reproduit ni des corrélations variant d'un
tenor à l'autre, ni la légère hausse des volatilités aux tenors extrêmes. Les estimations se lisent alors comme les
paramètres du modèle à un facteur le plus proche des moments observés. Les paramètres retenus pour la suite sont
\begin{equation}
a_x = 0{,}0107,\quad \sigma_x = 0{,}7367\,\%,\qquad a_y = 0{,}001,\quad \sigma_y = 0{,}3374\,\%,\quad \rho = 0{,}130.
\end{equation}

\subsection{Régimes de volatilité du spread}
\label{sec:regimes}

La méthode des moments suppose la volatilité du spread constante, alors que le test ARCH-LM rejette
l'homoscédasticité : $\sigma_y$ est une moyenne sur la période. Pour les spreads souverains de la zone euro,
\citet{MonfortRenne2014} emploient un modèle à deux régimes, dont un régime de crise
\citep[cours~5, slides~99-101]{MPR2025}. On ajuste le modèle de régression à changement de régime du cours
(slides~25 et~29) à la variation quotidienne du spread à vingt ans, expliquée par celle du swap €STR de même
tenor :
\begin{equation}
\Delta S_t = c_{k_t} + \beta_{k_t}\,\Delta R_t + \sigma_{k_t}\,\varepsilon_t,\qquad \varepsilon_t \sim \mathcal N(0,1)\ \text{i.i.d.},
\end{equation}
où le régime $k_t$, calme ou agité, suit une chaîne de Markov à deux états non observée. Le filtre de
Kitagawa-Hamilton donne la vraisemblance, maximisée en $(c_k, \beta_k, \sigma_k)$ et en les probabilités de
transition \citep[cours~5, slide~39]{Hamilton1989,MPR2025} ; l'algorithme de \citet{Kim1994} donne les
probabilités lissées des régimes. Dans chaque régime, la volatilité du spread et sa corrélation avec le swap se
calculent sur les variations pondérées par ces probabilités, et la part du temps passée dans chaque régime est sa
probabilité stationnaire (cours~5, slide~35).

\begin{table}[htbp]
\centering\small
\caption{Régimes de volatilité du spread à vingt ans (durées en jours ouvrés, volatilités en pb par an)}
\label{tab:regimes}
\begin{tabular}{@{}llccccc@{}}
\toprule
Variations & Régime & Part du temps & Durée moyenne & Vol. du spread & Vol. du swap & $\rho$ \\
\midrule
\multirow{2}{*}{Quotidiennes} & calme & 0,73 & 43 & 23,0 & 59,7 & 0,16 \\
 & agité & 0,27 & 16 & 51,1 & 89,4 & 0,15 \\
\midrule
\multirow{2}{*}{Hebdomadaires} & calme & 0,65 & 137 & 19,0 & 57,8 & 0,22 \\
 & agité & 0,35 & 73 & 42,6 & 95,2 & $-0{,}01$ \\
\bottomrule
\end{tabular}
\end{table}

Deux régimes se dégagent nettement (tableau~\ref{tab:regimes}). Sur variations quotidiennes, le régime calme
occupe 73~\% du temps, avec une volatilité du spread de 23~pb par an, et le régime agité 27~\%, avec 51~pb. Les
épisodes agités coïncident avec la remontée des taux de 2022 et du début de 2023, puis avec les crises politiques
françaises de juin 2024 (dissolution de l'Assemblée nationale) et de l'automne 2024 (budget, censure du
gouvernement) ; la figure~\ref{fig:regimes} les situe. Au 8 septembre 2026, la probabilité filtrée du régime
agité est de 8~\%.

\begin{figure}[htbp]
\centering
\includegraphics[width=\textwidth]{regimes}
\caption{Spread OAT–€STR à vingt ans et probabilité lissée du régime agité (variations quotidiennes).}
\label{fig:regimes}
\end{figure}

Deux résultats comptent pour la suite. La corrélation n'est pas liée aux régimes : presque la même dans les deux
régimes sur variations quotidiennes (0,16 et 0,15), elle s'annule dans le régime agité sur variations
hebdomadaires. Les régimes de volatilité n'expliquent donc pas la fourchette de $\rho$ selon les fenêtres. Les
régimes sont par ailleurs courts au regard des échéances d'options : 43 et 16 jours ouvrés en moyenne sur variations
quotidiennes, 137 et 73 sur variations hebdomadaires. La section~\ref{sec:resultats} en mesure l'effet sur le prix.
